\begin{textsample}{POS Dim 3 – rn – Score 16.00 – rn/rn\_inf\_7.txt}  \label{ex:f3_pos_057}
ESPECIFICIDADES DA \textbf{CRIANÇA} E CULTURAS \textbf{INFANTIS} Assume -se neste documento a \textbf{criança} como ser humano em fase inicial da vida – condição humana de inacabamento que confere abertura de possibilidades e potencial de desenvolvimento intenso, com semelhanças e especificidades em relação à indivíduos de outros ciclos de vida : vulnerabilidade-dependência do meio social ; capacidade para aprender e se desenvolver, de produzir cultura e participar da vida social ; integralidade de suas múltiplas dimensões humanas ( afetividade, cognição e motricidade ).
Com efeito, a infância deve sua diferença não à ausência de características ( presumidamente ) próprias do ser humano \textbf{adulto}, mas à presença de outras características distintivas que permitem que, para além de todas as distinções operadas pelo facto de pertencerem a diferentes classes sociais, ao género masculino e feminino, a seja qual for o espaço geográfico onde residem, à cultura de origem e etnia, todas as \textbf{crianças} do mundo tenham algo em comum ( SARMENTO, 2007, p. 35 ).
Sobre a vulnerabilidade/dependência da \textbf{criança} em relação aos \textbf{adultos} ( ou mais experientes ), para sua subsistência e desenvolvimento, Wallon ( 2005 ) destaca que a \textbf{criança} \textbf{pequena} necessita de \textbf{cuidados} físicos e psicológicos constantes.
Depende do outro para satisfação de suas necessidades durante um largo período de tempo, tanto nos aspectos físicos ( \textbf{equilíbrio}, locomoção, movimentos, alimentação, \textbf{higiene} e saúde, etc. ), como nos aspectos cognitivos ( construção de sua identidade e modos de exploração e significação do mundo/ da cultura ) e afetivos ( medo, raiva, choro, afeição, alegria, tristeza, etc. ).
Vigotski ( 2007 ) explica que a vulnerabilidade não se refere à incapacidade ou fraqueza da \textbf{criança}, mas à sua condição como pessoa que, estando ainda no início da vida, necessita da interação e \textbf{cuidado} de outros para que se desenvolva integralmente em todas as áreas. [... ] a imaturidade relativa da \textbf{criança}, em contraste com outras espécies, torna necessário um apoio prolongado por parte dos \textbf{adultos} [... ] por um lado ela depende totalmente de organismos imensamente mais experientes do que ela ; por outro lado, ela \textbf{colhe} os benefícios de um contexto ótimo e socialmente desenvolvido para o aprendizado ( VIGOTSKI, 2007, p. 166 ).
Nessa perspectiva, no cotidiano da Educação \textbf{Infantil}, compreender que ela é vulnerável, significa que, por exemplo, embora cuidando e protegendo -a para que não se machuque nas \textbf{brincadeiras}, para que não se queime nas horas de alimentar -se, para que não caia ao correr no \textbf{pátio} ou mesmo ao tentar subir numa árvore, ao mesmo tempo devemos possibilitar -lhe realizar tentativas e descobrir como se faz.
Devemos ainda contribuir para que aprenda e consiga, \textbf{progressivamente}, superar os desafios, sabendo os riscos que a rodeiam em cada ação que experimenta e encontrando, com \textbf{ajuda} do outro, as alternativas que a ajudarão na superação dos problemas.
Para tudo isso é necessário que a \textbf{criança} sinta -se segura, sabendo que haverá alguém que poderá ajudá -la quando precisar.
Uma característica que diferencia/especifica a \textbf{criança} em relação aos \textbf{adultos} é a globalidade, que reflete a forma holística pela qual ela aprende e se desenvolve.
Para Zabalza ( 1998 ; 2008 ) a \textbf{criança} é um sujeito não setorizável.
É inteira, se desenvolve como um todo integrado no qual o afetivo, o social e o cognitivo se constituem numa dinâmica intensa, não sendo possível, portanto, tratar isoladamente, no processo educativo, apenas uma dessas “ partes ”, priorizando um aspecto em detrimento dos demais.
Por fim, Oliveira-Formosinho ( 2005 ) destaca que a Psicologia tem chamado a atenção, nas últimas duas décadas, para o fato de que a competência da \textbf{criança}, desde muito \textbf{pequena}, tem uma latitude mais acentuada do que a sua vulnerabilidade social aparenta.
Ou seja, reconhece na \textbf{criança} humana, suas competências sociopsicológicas que se manifestam desde a mais tenra idade, por exemplo, nas suas formas precoces de comunicação – em cem linguagens.
Neste sentido, compreende -se que ao mesmo tempo em que a \textbf{criança} é vulnerável/dependente do \textbf{adulto}, ela é capaz de aprender e se desenvolver, desde que lhe sejam dadas as condições propícias.
Todas as \textbf{crianças} têm capacidade/competência para aprender, salvo aquelas com alguma deficiência ou transtorno fora da normalidade.
No entanto, o quê e como elas aprendem varia conforme as práticas culturais que vivenciam.
Desta forma, é relevante o papel das interações com os objetos, o meio social e os signos, e a mediação do outro nas práticas culturais.
Sendo assim, tanto no âmbito da Psicologia como da Sociologia, destaca -se que outra importante especificidade \textbf{infantil} é a capacidade de produzir cultura.
Nas suas interações sociais, entre pares, as \textbf{crianças} produzem culturas \textbf{infantis} relativas à fantasia, à imaginação e à \textbf{brincadeira}.
Segundo Kramer ( 2006, p. 16 ) “ a cultura \textbf{infantil} é, pois, produção e criação.
As \textbf{crianças} produzem cultura e são produzidas na cultura em que se inserem ( em seu espaço ) e que lhes é contemporânea ( de seu tempo ) ”.
Sarmento ( 2007 ) afirma que : “ todas as \textbf{crianças}, desde \textbf{bebês}, têm múltiplas linguagens ( gestuais, corporais, \textbf{plásticas} e verbais ) por que se expressam ” ; têm “ outras racionalidades ”, construídas “ nas interacções de \textbf{crianças}, com a incorporação de \textbf{afetos}, da fantasia e da vinculação ao real ” ; [... ] “ todas as \textbf{crianças} trabalham, nas múltiplas tarefas que preenchem seus cotidianos, na escola, no espaço \textbf{doméstico} e, para muitas, também nos campos, nas oficinas ou na rua ” ( SARMENTO, 2007, p. 36 ).
Finalmente, segundo o autor, “ [... ] as \textbf{crianças} são “ sujeitos activos, que interpretam e agem no mundo.
Nessa acção, estruturam e estabelecem padrões culturais.
As culturas \textbf{infantis} constituem, com efeito, o mais importante aspecto na diferenciação da infância ” ( Ibid, p. 36 ).
Para que se compreenda o que significa cultura da infância, é necessário entender o que se assume como cultura.
No campo da Antropologia, a terminologia “ cultura ” corresponde aos diversos modos de vida, valores e significados compartilhados por diferentes grupos e períodos históricos.
Apresenta -se também, em práticas culturais, atividades coletivas que variam de grupo social para grupo social.
É uma prática social entendida não como um fazer ( artes ) ou estado de ser ( civilização ), mas como conjunto de práticas por meio das quais significados são produzidos e compartilhados em um grupo ( MOREIRA ; CANDAU, 2007 ).
Quando um grupo compartilha uma cultura, socializa um conjunto de significados construídos, ensinados e aprendidos por meio de suas linguagens.
É assim que também acontece na infância.
A cultura é firmada a partir do que é produzido por e para as \textbf{crianças} através das produções simbólicas que estão associadas às interações, às \textbf{brincadeiras}, à fantasia, etc.
Assim, a capacidade que a \textbf{criança} tem de significar o mundo e constituir seus modos próprios de ser e de agir, contribuindo com a transformação dos espaços que ocupa, configura -se em cultura \textbf{infantil}.
Nesse sentido, as \textbf{crianças} produzem cultura, pois são capazes de \textbf{inventar}, ampliar os modos de se relacionar com o mundo e compreender como este funciona.
A Cultura da Infância é marcada pela ludicidade, produzida nas interações entre pares e intergeracionais, nas \textbf{brincadeiras} e linguagens \textbf{infantis}.
Caracteriza -se por seu caráter coletivo – é no interior do grupo de \textbf{crianças} que estas coletivamente produzem uma cultura singular.
Os produtos da cultura \textbf{infantil} expressam esse caráter coletivo de participação.
A materialidade do produto – fruto das interações coletivas – é tão importante quanto o processo de produção.

% matched lemmas: adulto, afeto, ajuda, bebê, brincadeira, colher, criança, cuidado, doméstico, equilíbrio, higiene, infantil, inventar, pequeno, plástico, progressivamente, pátio
\end{textsample}
